\documentclass[12pt,a4paper]{report}

% 1. ตั้งค่าระยะขอบกระดาษ (ซ้าย 1.5 นิ้ว, ขวา 1 นิ้ว, บน 1 นิ้ว, ล่าง 1 นิ้ว)
\usepackage[left=1.5in, right=1in, top=1in, bottom=1in]{geometry}

% 2. ตั้งค่าระบบภาษาไทย และ XeLaTeX
\usepackage{fontspec}
\usepackage{polyglossia}
\setdefaultlanguage{thai}
\setotherlanguage{english}
\XeTeXlinebreaklocale "th"
\XeTeXlinebreakskip = 0pt plus 1pt

% 3. นำเข้าฟอนต์ TH Sarabun New
\setmainfont[
    Path=fonts/,
    Extension=.ttf,
    BoldFont=THSarabunNew_Bold,
    ItalicFont=THSarabunNew_Italic,
    BoldItalicFont=THSarabunNew_BoldItalic,
]{THSarabunNew}

\usepackage{setspace}

\begin{document}

% -------------------------
% หน้าปกนอก (Cover Page)
% -------------------------
\begin{titlepage}
\begin{center}
    \setstretch{1.5}
    % ชื่อเรื่องภาษาไทย
    {\fontsize{20pt}{24pt}\selectfont \textbf{ชื่อเรื่องวิทยานิพนธ์ภาษาไทย}}\\
    \vspace{0.5cm}
    % ชื่อเรื่องภาษาอังกฤษ
    {\fontsize{20pt}{24pt}\selectfont \textbf{THESIS TITLE IN ENGLISH}}\\
    
    \vspace{2.5cm}
    
    % ชื่อผู้จัดทำ
    {\fontsize{18pt}{22pt}\selectfont \textbf{ชื่อ นามสกุล}}\\
    {\fontsize{18pt}{22pt}\selectfont \textbf{NAME SURNAME}}\\
    
    \vfill
    \setstretch{1.15}
    
    % ส่วนท้ายของหน้าปก
    {\fontsize{18pt}{22pt}\selectfont \textbf{วิทยานิพนธ์นี้เป็นส่วนหนึ่งของการศึกษาตามหลักสูตรปริญญาวิทยาศาสตรมหาบัณฑิต}}\\
    {\fontsize{18pt}{22pt}\selectfont \textbf{สาขาวิชาเทคโนโลยีสารสนเทศ}}\\
    {\fontsize{18pt}{22pt}\selectfont \textbf{คณะเทคโนโลยีสารสนเทศ}}\\
    {\fontsize{18pt}{22pt}\selectfont \textbf{สถาบันเทคโนโลยีพระจอมเกล้าเจ้าคุณทหารลาดกระบัง}}\\
    {\fontsize{18pt}{22pt}\selectfont \textbf{พ.ศ. 25XX}}\\
    
    \vspace{1.5cm}
    
    % รหัสวิทยานิพนธ์
    {\fontsize{16pt}{20pt}\selectfont \textbf{KMITL-20XX-IT-M-XXX-XXX}}
\end{center}
\end{titlepage}

% -------------------------
% หน้าปกใน (ภาษาอังกฤษตามฟอร์แมต)
% -------------------------
\begin{titlepage}
\begin{center}
    \setstretch{1.5}
    {\fontsize{20pt}{24pt}\selectfont \textbf{THESIS TITLE IN ENGLISH}}\\
    
    \vspace{2.5cm}
    
    {\fontsize{18pt}{22pt}\selectfont \textbf{NAME SURNAME}}\\
    
    \vfill
    \setstretch{1.15}
    
    {\fontsize{18pt}{22pt}\selectfont \textbf{A THESIS SUBMITTED IN FULFILLMENT}}\\
    {\fontsize{18pt}{22pt}\selectfont \textbf{OF THE REQUIREMENT FOR THE DEGREE OF}}\\
    {\fontsize{18pt}{22pt}\selectfont \textbf{MASTER OF SCIENCE IN INFORMATION TECHNOLOGY}}\\
    {\fontsize{18pt}{22pt}\selectfont \textbf{SCHOOL OF INFORMATION TECHNOLOGY}}\\
    {\fontsize{18pt}{22pt}\selectfont \textbf{KING MONGKUT'S INSTITUTE OF TECHNOLOGY LADKRABANG}}\\
    {\fontsize{18pt}{22pt}\selectfont \textbf{20XX}}\\
    
    \vspace{1.5cm}
    {\fontsize{16pt}{20pt}\selectfont \textbf{KMITL-20XX-IT-M-XXX-XXX}}
\end{center}
\end{titlepage}

% -------------------------
% หน้าลิขสิทธิ์
% -------------------------
\clearpage
\vspace*{5cm}
\begin{center}
    {\fontsize{18pt}{22pt}\selectfont \textbf{COPYRIGHT 20XX}}\\
    \vspace{0.5cm}
    {\fontsize{18pt}{22pt}\selectfont \textbf{SCHOOL OF INFORMATION TECHNOLOGY}}\\
    \vspace{0.5cm}
    {\fontsize{18pt}{22pt}\selectfont \textbf{KING MONGKUT'S INSTITUTE OF TECHNOLOGY LADKRABANG}}
\end{center}
\clearpage

% -------------------------
% บทคัดย่อภาษาไทย
% -------------------------
\chapter*{บทคัดย่อ}
\addcontentsline{toc}{chapter}{บทคัดย่อ}
พิมพ์เนื้อหาบทคัดย่อภาษาไทยที่นี่... สามารถพิมพ์และตัดคำภาษาไทยได้สมบูรณ์ครับ

% -------------------------
% บทคัดย่อภาษาอังกฤษ (ABSTRACT)
% -------------------------
\chapter*{ABSTRACT}
\addcontentsline{toc}{chapter}{ABSTRACT}
Write your abstract in English here...

% -------------------------
% สารบัญ
% -------------------------
\tableofcontents
\clearpage

% -------------------------
% บทที่ 1
% -------------------------
\chapter{บทนำ}
ส่วนเนื้อหาของบทที่ 1 เริ่มต้นตรงนี้...

\end{document}